% New composite audit; all component proofs remain unchanged.
\section*{Independent assembly audit of the restricted retained-phase convolution estimate}
\label{sec:restricted-prop52-composite-20261009}

The conclusion checked here is a one-bounded convolution statement.
It is not yet a prime-product distribution theorem, a proof of the
original unrestricted Proposition 5.2, or a restoration of $K,c,b$.
The component spectral theorem and the classical Shiu, Weil, Poisson
and SW inputs are cited inputs, not re-proved in this assembly audit.

\paragraph{The statement actually obtained.}
Fix $c_0>0$, $0<\sigma<1/100$, and
$0\leq\kappa\leq\sigma/100$.  Set $x=NM$, $\theta=7/32$, and assume
\[
 c_0x\leq |a|\leq x^{1+\kappa},\qquad
 Q_1<Nx^{-\sigma},\qquad N^2Q_2Q_3^2<x^{1-8\sigma},
\]
\[
 |a|^\theta NQ_1Q_2^5Q_3^2<x^{2-15\sigma},\qquad
 N^{1+2\theta}Q_1Q_2^{5-2\theta}Q_3^{2+2\theta}<x^{2-15\sigma}.
\]
All five sequences $\alpha_n,\beta_m,\gamma_{q_1},\lambda_{q_2},
\nu_{q_3}$ are one-bounded.  The modulus sequences have support
$(q_i,a)=1$.  The sequence $\alpha$ is supported on
$P^-(n)>z_0=x^{1/(\log\log x)^3}$ and satisfies the full auxiliary
SW condition: for every $A>0$, uniformly in $(b,q)=1$ and $d\geq1$,
\[
 \sum_{\substack{n\sim N\\n\equiv b\ (q)\\(n,d)=1}}\alpha_n
 -{1\over\varphi(q)}\sum_{\substack{n\sim N\\(n,dq)=1}}\alpha_n
 \ll_A N\tau(d)^{O(1)}(\log N)^{-A}.
\]
If the sequences depend on the residue, these constants must be
uniform in that family.  Define the discrepancy $\Delta_a(l)$ with
the original coprime mean, exactly as L23 line 2448.  Then the assembled
result is, for every fixed $A>0$,
\begin{equation}\label{eq:restricted-composite-conclusion}
 \left|\sum_{q_i\sim Q_i}\gamma_{q_1}\lambda_{q_2}\nu_{q_3}
                 \Delta_a(q_1q_2q_3)\right|
 \ll_{A,\sigma,c_0,\mathrm{SW}} x(\log x)^{-A}.
\end{equation}
Uniformity is over residues satisfying their own displayed
$|a|^\theta$ condition; the supremum does not erase that condition.
The derivation below establishes this restricted statement relative
to the component proofs' explicit external inputs.

\paragraph{D: the preliminary reduction is applicable.}
Writing $n,q,u,v$ for the logarithms to base $x$ of $N,Q_1,Q_2,Q_3$,
the four conditions imply
\[
 N>x^\sigma,\quad N<x^{1/2-4\sigma},\quad
 Q_1Q_2Q_3<x^{7/10-13\sigma/2}.
\]
The last inequality follows from $q<n-\sigma$,
$n+u/2+v<(1-8\sigma)/2$, and $u<(2-15\sigma)/5$.
With $d\sim D\leq y_0^2=x^{o(1)}$ and
$D\mathcal R\asymp Q_2Q_3$, one has, eventually,
\[
 N,M\geq x^{\delta},\quad Q_1D\leq Nx^{-\delta},\quad
 Q_1D\mathcal R\leq x^{1-\delta},\qquad\delta=\sigma/3.
\]
Every scale and the residue are polynomially bounded.  These are
exactly the additional margins of the zero-product supplement.
Its shifted moments omit $mn=a$ before the variance is expanded;
restoration occurs only after the dangerous double exact diagonal
has been excluded.  The normalized restoration error is
$x^{1-\delta/2+o(1)}$, not a factor times a merely logarithmic error.

\paragraph{The small factors do not multiply a logarithmic remainder.}
The weighted smooth-part deletion is performed first on the grouped
modulus $l=q_1q_2q_3<x^{1-\delta}$.  Primitive Shiu estimates and
Rankin's bound handle the divisor multiplicity inside the sum.
For the retained terms write $q_2=us_2,q_3=vs_3$, and
$d=s_2s_3$, $r=uv$.  The coefficient
\[
 \Lambda_{d,r}=\sum_{s_2s_3=d}\sum_{uv=r}
              \lambda_{us_2}\nu_{vs_3},\mathbf1_{\rm fixed\ supports}
\]
is independent of $q_1$ and bounded by $\tau(d)\tau(r)$.
The whole dyadic $d$ range is supplied to D once.  The original
$\gamma_{q_1}$ is its outer bounded coefficient, and disappears in
the indicated Cauchy bound; it is not an unaccounted extra factor
in either exponential sum.  Only logarithmically many dyadic ranges
are added.  In particular, no $x^{o(1)}$ separate values of $d$
multiply D's $x\log^{-A}x$ remainder.

\paragraph{E1: the near-linear lower bound on $|a|$ is essential.}
The checked replacement uses Maynard I v2 Lemma 17.2 directly.
Since $(a,qr_1r'_1)=1$, the completed Weil gcd loses its $a$ factor
exactly.  The algebraic consequence of the size conditions is
\[
 n+\tfrac52q+3u+3v
 <{1021\over512}-{147\over8}\sigma+O(1/\log x).
\]
It follows by the nonnegative combination $35/16,23/16,5/16$
of $q-n$, $2n+u+2v$, $n+q+5u+2v$, and subtracting $v/2$.
The third bound uses $|a|\geq c_0x$ and $\theta=7/32$;
neither can be removed here.  The other Weil conditions follow from
$n+q<1-9\sigma$ and $2n+q+u+v<17/10-29\sigma/2$.

For E1 alone it is legitimate to set $q'=dq$, because that lemma
permits arbitrary bounded coefficients depending on $q'$ and $r$.
The sparse indicator $\mathbf1_{d\mid q'}$, and the square root of
the nonnegative smooth $q$ weight, fit this permission.  Its actual
denominator is $\varphi(q'r_2)q'r_1$.  For each fixed $d$ the bound
is $N^2/(dQ_1)x^{-\sigma+o(1)}$.  Summing $1/d$ and spending the
remaining subpower $D$ gives
\begin{equation}\label{eq:restricted-E1}
 |E_1|\ll {N^2\over DQ_1}x^{-\sigma/2}.
\end{equation}

\paragraph{RP and S: three distinct polarizations, no arbitrary $q$ weight.}
The second error has the exact identity
\[
 E_2(a)=\sum_{d\sim D}{E_3(a;d)\over d}.
\]
The denominator in $E_3$ is $qr_1r_2$, without an additional $d$.
Reciprocity retains $e(ah/(qdr_1r_2n_1))$.  Its normalized coefficient
is $(a/x)(D/d)\log^5x$, so the five-variable variation costs
$x^{5\kappa+o(1)}$ uniformly in $d\sim D$.
All endpoints from summation by parts are frozen before S is used.
The cross-coprime partition applies to the rough pairs
$(n_1r_2,n_2)$ and costs only $x^{o(1)}$.

Expand the two factors of $\Lambda$ independently at this stage.
The choices of smooth factors on the two sides need not agree.
There are $\tau(d)^2$ choices.  The two resulting rough-factor
coefficient pairs (called $\Gamma$ and $\Lambda$ in S), and the
independently truncated $n_1$ coefficient paired with $\alpha_{n_2}$,
are three separate sesquilinear pairs.  For each pair use
\[
 B(u,v)=\tfrac14\sum_{j=0}^3i^jB(u+i^jv,u+i^jv).
\]
After normalizing the vectors by $2$, this costs at most $4^3=64$.
The kernel need not be Hermitian.  SW was used before this operation;
S requires only bounded coefficients.  The product endpoints on
the rough factors and the independent fixed $h$ weights are the
precise extensions previously checked, not arbitrary coupled
arithmetic weights.

Keep the original $q$ and fixed $d$ in S.  Its original statement
already has congruence modulo $qd$.  Unlike E1, it is not licensed
to replace $q$ by $qd$ and impose a sparse $q$ indicator.
With $S=Q'_2$, $R=Q'_3$, all monomial powers are positive, including
$5-2\theta$.  The $8\sigma,15\sigma$ hypotheses therefore imply
S's $7\sigma,14\sigma$ conditions after all fixed dyadic constants.
The Fourier length has the allowed factor $D\log^5x=x^{o(1)}$.

At the actual spectral substitution $z=s_1s_2$, both $s_i$ still
satisfy $(s_i,a)=1$.  Thus $(a,z)=1$, as required by the retained
packet proof.  The packet norm preserves the interval coefficient
on the other spectral side; no unrestricted non-coprime version
of the printed spectral theorem has been inserted.
Taking just the deliberately weaker pre-variation bound from S,
\[
 |E_3(a;d)|\ll {N^2\over Q_1}x^{-\sigma+5\kappa+o(1)}.
\]
The sum of $1/d$ is $O(1)$.  Allocate a total $\sigma/10$ to
the remaining subpower factors, including conversion to $1/D$.
Since $5\kappa\leq\sigma/20$, at least $17\sigma/20$ remains;
weakening to $\sigma/2$ yields
\begin{equation}\label{eq:restricted-E2}
 |E_2|\ll {N^2\over DQ_1}x^{-\sigma/2}.
\end{equation}
This conservative accounting does not need to spend all of S's
stronger natural-scale reserve.

\paragraph{Normalization and the actual final error.}
Substitution of \eqref{eq:restricted-E1}--\eqref{eq:restricted-E2}
into D gives
\[
 M(DQ_1)^{1/2}\bigl(|E_1|^{1/2}+|E_2|^{1/2}\bigr)
 \ll x^{1-\sigma/4},
\]
before fixed logarithmic factors.  Separately D's zero-product
correction is $x^{1-\sigma/6+o(1)}$.  This is weaker than the
exponential-sum saving and must not be called $x^{1-\sigma/5}$.
Both terms are absorbed by $x\log^{-A}x$ for each fixed $A$.
D already has an arbitrary-log SW remainder, so no overall
fixed-power saving is claimed.  Increase its logarithmic exponent
before summing the finitely many dyadic decompositions; this proves
\eqref{eq:restricted-composite-conclusion}.

\paragraph{Dependence and limits of the audit.}
The order is D and its independent zero-product review; E1 and its
independent Weil review; the retained-phase elementary reduction;
the coprime spectral audit S and the independent packet/scale review;
then this composition.  None uses the conclusion being assembled.
This author supplied S, so this report does not describe S as a new
independent reproof; the root packet review provides its separate
cross-check.  The other component proofs and the bridge were read
directly here, not accepted from their status labels.

This assembly supplies the restricted one-bounded statement above.
It does not supply divisor-bounded HB coefficient truncation,
one/two-long decompositions, prime products at all local scales,
small-product tails, finite sieve compilation, or the reported
unidentified Wu correction.  Those downstream matters cannot be
inferred merely by changing the status of the present statement.
